\chapter{Characterisation of Clusters}\label{ch:candidate_characterisation}

\section{Introduction}\label{sec:ch6_introduction}

Chapter~\ref{ch:cluster_transition_screening} identified candidate superspreading-compatible (CSC) clusters from their size, membership novelty, and amount of later linked sequences. This chapter asks whether those clusters represent sampled infections concentrated within particular populations, infections distributed across populations, or different combinations of these patterns. The purpose is to develop epidemiological questions for follow-up investigation and to establish how much the profiles selected by the screen differ from the sampled background.

There are reasons to examine both who is represented and how broadly a cluster spans population groups. Social contact surveys show that contacts are strongly structured by age, especially among schoolchildren and young adults \citep{Mossong2008}. Scottish genomic investigations identified community, healthcare, and conference outbreaks following viral introductions \citep{daSilvaFilipe2020}; the EURO 2020 investigation linked a distinctive demographic profile to reported football-related activities \citep{Marsh2021}. Outbreaks can also differ in their subsequent reach: a Boston nursing-facility outbreak involved intense local transmission with little wider spread, whereas a conference outbreak was followed by extensive dissemination \citep{Lemieux2021}. These studies motivate questions about population concentration and geographic breadth, although residential and demographic attributes cannot identify exposure settings by themselves.

\emph{Composition} identifies the population groups represented among sampled infections. \emph{Within-cluster diversity} measures how evenly those records are distributed across the categories of an attribute. Shannon entropy provides a compact summary of the latter. A cluster whose members mostly belong to one age group raises different questions about shared activities from a cluster with a more even age distribution. A cluster spanning health boards raises questions about connections between communities. Entropy is useful here as a description of these patterns, rather than as a direct measure of social contact or transmission. Category proportions remain necessary because entropy does not identify which groups predominate.

The contribution of this chapter is to apply these descriptions systematically after a screen that did not use the demographic and geographic attributes to assign CSC status. This complements the geographic structure described in Chapter~\ref{ch:genomic_networks} and the screening results in Chapter~\ref{ch:cluster_transition_screening}. It does not validate CSC status against confirmed superspreading events. We consider two exploratory explanations: demographic or regional concentration may accompany shared settings and local bursts, whereas broader geographic representation may accompany connections associated with onward burden. These are hypotheses for subsequent epidemiological investigation, not hypotheses specified before inspecting these data. We do not assume that every attribute, including deprivation, should have a monotonic association with screening outcomes.

We first assess the size distributions of CSC and background clusters, then compare composition and diversity across cluster sizes and screening routes. Diversity is described both absolutely and relative to random samples of the same size as the cluster, using category proportions from the same surveillance window. Bayesian hierarchical models examine associations with CSC status and the two screening scores, accounting for other measured attributes, policy period, and viral clade, with additional surveillance adjustment in expanded models. All analyses concern the sampled record; they do not estimate individual transmission risk or establish causes.

\section{Materials and Methods}\label{sec:ch6_methods}

We used the transition graph and screening outcomes from Chapter~\ref{ch:cluster_transition_screening}. The screen started with 277,154 clusters in retained surveillance windows before applying the size threshold. Of these, 8,962 (3.2\%) met the primary minimum-size threshold of six sequences and formed the characterisation dataset. We assigned 631 clusters (7.0\% of size-eligible clusters) a candidate superspreading-compatible (CSC) status at calibrated $p \leq 0.05$, based on high-priority classifications for local bursts, onward burden, or both. All other clusters meeting the size threshold form the background group, including those labelled possible review at $0.05<p\leq0.10$ on an applicable screening axis.

We analysed three cluster-level outcomes: binary CSC status (candidate or background), local burst score, and onward burden score. Each score is the mean of its available components' percentile ranks within the same retained surveillance window and lies on a 0--1 scale. Onward burden scores are available for only 4,216 of the clusters meeting the size threshold: those with positive observed downstream burden, referred to as burden-eligible clusters. We characterised screening outcomes at two analytical levels summarised in Table~\ref{tab:ch6_analytical_levels}: \emph{cluster-level} analyses use one row per cluster to quantify within-cluster diversity, while \emph{sequence-level} analyses use one row per sequence record to describe the distribution of categories among sampled infections.

\begin{table}[htbp]
	\centering
	\caption[Analytical levels for characterisation of screened clusters]{Analytical levels for characterisation of screened clusters. Cluster-level analyses examine within-cluster diversity; sequence-level analyses examine sampled population composition.}\label{tab:ch6_analytical_levels}
	\begin{thesistablebody}{@{}P{0.18\linewidth}P{0.38\linewidth}P{0.38\linewidth}@{}}
		\toprule
		\textbf{Feature} & \textbf{Cluster-level} & \textbf{Sequence-level} \\
		\midrule
		Unit of analysis & One row per cluster in a retained window & One row per sequence record \\
		Outcomes & CSC status and local burst score; onward burden score for burden-eligible clusters & Corresponding cluster outcomes; onward burden score only for records in burden-eligible clusters \\
		Predictors & Five attribute diversity measures (sex, age, SIMD, urban/rural class, health board) & Five categorical attributes (sex, age group, SIMD quintile, urban/rural class, health board) \\
        Purpose & Relate outcomes to within-cluster diversity & Relate outcomes to sampled population composition \\
		\bottomrule
	\end{thesistablebody}
\end{table}

The following sections define the attributes, entropy measures, descriptive comparisons, and Bayesian models. Appendix~\ref{app:candidate_characterisation} provides model details, diagnostic guidance, and supplementary results.

\begin{tcolorbox}[
    breakable,
    title=\textbf{Glossary}: key terms used in this chapter,
    halign=right, 
    fontupper=\small
    ]
\begin{description}
    \item[Sequence record:] A record of a sequenced infection in a retained surveillance window, assigned to a cluster. Each record includes sex, age group, SIMD quintile, urban/rural class, health board, and measures of the epidemic context at sampling. Each record is one observation in sequence-level analyses.
    \item[Cluster-level analysis:] One observation per cluster in a retained window; used to relate diversity among member records to the cluster's screening outcomes.
    \item[Sequence-level analysis:] One observation per sequence in each retained window in which it appears; used to relate population categories to the screening outcomes of the containing clusters.
    \item[Composition:] The proportions of observed records in each category of sex, age group, SIMD quintile, urban/rural class, or health board attribute.
    \item[Within-cluster diversity:] How evenly records are distributed across the categories of one attribute. The model tables use \emph{mixing} as a shorthand; it describes sampled attributes, not observed contacts.
    \item[Observed entropy:] A summary of within-cluster diversity on a 0--1 scale: zero means that all observed records share one category; one means equal representation of every category defined for that attribute.
    \item[Null-standardised entropy:] Observed entropy minus its expected value under random category assignment from the same window, divided by the simulated standard deviation. The reference also matches the cluster's size. Negative values indicate less diversity than expected; positive values indicate more. This scale is not bounded by 0 and 1.
    \item[CSC status:] A yes/no outcome indicating whether a cluster meeting the size threshold received candidate superspreading-compatible (CSC) status under the primary screening criteria (see Chapter~\ref{ch:cluster_transition_screening} for details).
    \item[Local burst score:] A summary of unusually large accumulation of sampled infections and, where earlier links exist, new cluster membership. Here, local refers to a cluster and its immediate graph links, not necessarily one geographic area. The score averages within-window percentile ranks for contextual excess size and available upstream novelty (see Chapter~\ref{ch:cluster_transition_screening} for details).
    \item[Onward burden score:] A summary of additional sampled infections in later linked clusters relative to the starting cluster's size. It averages percentile ranks for direct and cumulative burden, and is available only when observed downstream burden is positive. Graph links record shared membership, not established transmission (see Chapter~\ref{ch:cluster_transition_screening} for details).
    \item[Policy period:] A phase of Scottish COVID-19 policy used to represent changes in epidemic and surveillance context, including opportunities for contact (see Chapter~\ref{ch:genomic_networks} for details).
    \item[Nextclade clade:] A broader grouping of related viral lineages \citep{Aksamentov2021}. Clade intercepts represent differences between viral groups and their associated epidemic context; they do not isolate biological effects.
\end{description}
\end{tcolorbox}

\subsection{Analysis variables}\label{sec:ch6_attributes}

We used five sequence attributes from Chapter~\ref{ch:genomic_networks}: sex, age group, population-weighted Scottish Index of Multiple Deprivation (SIMD) quintile, urban/rural class, and health board. These enter the cluster-level regressions as continuous measures of within-cluster diversity and the sequence-level regressions as categorical predictors. Both analytical levels use logistic regression for CSC status and linear regression for the two screening scores. Local authority was additionally used in the descriptive entropy comparisons. All six attributes were recorded for every sequence record. Reference groups for the sequence-level categorical predictors are male, age 25--64, SIMD quintile 1 (most deprived), large urban areas, and Greater Glasgow and Clyde.

Policy period and Nextclade clade describe the timing of sampling and the viral group. Each sequence's policy period is assigned from its collection date; the cluster-level value is the most common policy period among its records. Each cluster contains sequences from a single lineage, whereas a clade may encompass multiple lineages. Consequently, sequences from different lineages can belong to the same clade, a hierarchical relationship utilised at both analytical levels. Policy period and clade have separate varying intercepts, allowing the model baseline to differ between periods and between clades. In Chapter~\ref{ch:genomic_networks}, policy period and clade explained most variation in the largely geographic patterns of compatibility scores between pairs of infections.

In addition to policy period and clade, we used three surveillance measures to assess whether the findings change after accounting for differences in testing and sequencing intensity:

\begin{description}
    \item[Window sequencing proportion:] Unique sequences divided by positive tests within the surveillance window.
    \item[Datazone cumulative incidence:] Positive tests accumulated from the start of the available testing series (1 July 2020) up to and including the collection date, divided by the datazone population. This describes accumulated recorded epidemic burden, not incidence within the current surveillance window.
    \item[Datazone cumulative sequencing proportion:] \mbox{}\\
    Unique sequences accumulated from the available quality-filtered metadata (starting 4 July 2020) divided by positive tests accumulated from the available testing series, both counted up to and including the collection date. This describes local sequencing coverage of recorded positive tests up to sampling; it is not restricted to the current window.
\end{description}

Positive tests include PCR and lateral-flow results. Sequence-level analyses use non-missing values for each collection date; cluster-level analyses average these values across records with data. Cumulative sequencing proportion is undefined when cumulative recorded positive tests are zero; cumulative incidence is zero in that case. These cumulative measures cannot recover infections before the available surveillance series or undetected infections. Window sequencing proportion is shared by all sequences and clusters in the same window.

We standardised the three surveillance measures separately at cluster and sequence level by subtracting the mean and dividing by the sample standard deviation (SD)\nomenclature{SD}{standard deviation}. These calculations used all non-missing values across retained windows, before applying the size threshold or excluding records with missing data required by each model. Each value contributed equally, so windows with more contributing records had greater weight. A one-unit increase therefore represents one SD in the corresponding dataset before filtering (Appendix~\ref{sec:app_ch6_bayesian_details}).

\subsection{Within-cluster diversity measures}\label{sec:ch6_composition_entropy}

We summarised how evenly the observed records in each cluster were distributed across the categories of an attribute. These categories include demographics (for example, age group) and residential characteristics (for example, health board). Entropy combines the number of represented categories with their relative frequencies. For a two-category attribute, two clusters of ten records with splits of 10:0 and 5:5 have normalised entropies of 0 and 1, respectively. Splits of 9:1 and 1:9 have the same entropy, despite opposite composition. Thus entropy can identify concentration or breadth, while category proportions are needed to identify which population is represented. We use \emph{observed entropy} for absolute diversity and \emph{null-standardised entropy} for diversity relative to a window and size-specific reference.

\subsubsection*{Observed entropy}

For categorical attribute $v$ (e.g., age group) in cluster $i$, let $p_{ivk}$ be the proportion of records in category $k$ (e.g., 15--24), and $K_v$ the number of distinct categories observed across all retained windows before size restriction (e.g., $K_v=6$ for age group). The observed normalised Shannon entropy $H_{iv}^{\mathrm{obs}}$ is:

\[
    H_{iv}^{\mathrm{obs}} = -\frac{1}{\log_2 K_v}\sum_{k=1}^{K_v} p_{ivk}\log_2(p_{ivk}),
\]

adopting $0\log_2 0=0$ (the standard convention that empty categories contribute zero to entropy). This quantifies category diversity on a 0--1 scale: entropy of 0 indicates all sequences in the cluster share one category; entropy of 1 indicates all $K_v$ categories are equally represented. Normalisation gives each attribute a common scale, though the diversity a cluster can reach depends on its size and the available categories.

\subsubsection*{Null-standardised entropy}

Infection risk, vaccination eligibility, and transmission and surveillance intensity varied over time and across population groups (Chapter~\ref{ch:genomic_networks}). These factors affect cluster size and composition, which in turn affect observed entropy. We therefore constructed a simulated reference for each attribute to account for differences in cluster size and window composition.

For each attribute, we pooled all sequence records in each retained window, including records in CSC, background, and size-ineligible clusters. Each record contributed once to that window's category proportions. For each distinct cluster size of at least two in that window, we simulated 1,000 groups of that size by drawing categories with replacement according to those proportions (multinomial sampling). We calculated entropy for each simulated group. Clusters of the same size in the same window shared the reference distribution for that attribute, regardless of CSC status.

This reference represents the entropy expected under random category assignment given the window's composition and the cluster's size; it does not account for lineage structure or preserve associations between attributes. Null-standardised entropy is:

\[
    Z_{iv}^{(H)} = \frac{H_{iv}^{\mathrm{obs}} - \mu_{iv}^{\mathrm{null}}}{\sigma_{iv}^{\mathrm{null}}},
\]

where $\mu_{iv}^{\mathrm{null}}$ and $\sigma_{iv}^{\mathrm{null}}$ are the simulated mean and sample standard deviation. A value of 0 means diversity matches expectation; positive values indicate greater diversity; negative values indicate less. Observed entropy is zero for a cluster containing one sequence. The null-standardised value is undefined if the simulated reference is unavailable or has no variation. Standardisation supplies a size-specific reference; it does not guarantee that the resulting values are independent of cluster size. The reference standard deviation itself changes with size, so the same absolute departure from expected entropy can yield different standardised values.

\subsection{Descriptive comparisons}\label{sec:ch6_descriptive_methods}

We first compared the distribution of cluster sizes between CSC and background groups. Cumulative plots show both the number of clusters of size at most $s$ and the sum of their record counts, with corresponding percentages within each group. A sequence appearing in two retained windows contributes a record in each window; these totals are not counts of unique people or unique sequences across the study.

We then used five size bands: 6--9, 10--19, 20--49, 50--99, and $\geq100$ sequences. Within each band, the main composition summaries pool records within CSC and background groups and divide category counts by the total number of records in the corresponding group. Larger clusters therefore contribute more records. We report both group percentages and their difference in percentage points (CSC minus background). Appendix comparisons average within-cluster category proportions instead, giving each cluster equal weight. Tables report the numbers of contributing clusters and records.

We describe the distribution of each entropy measure using its median and interquartile range (IQR), separately by size band and CSC status. Each cluster receives equal weight. The IQR describes variation between clusters, not uncertainty in an estimated median.

We also summarised burst-only, burden-only, and both-axis candidates separately. Burst-only candidates were compared with all eligible background clusters; burden-only and both-axis candidates were compared with the burden-eligible subset of the background. The latter comparison respects the requirement for positive observed downstream burden. The small both-axis group was retained for description rather than separate regression.

Finally, we assessed whether the descriptive profiles depended on the screening cutoff or on including smaller clusters. We reapplied $p\leq0.01$, $0.025$, $0.05$, and $0.10$ to the saved calibrated scores, retaining the six-sequence testing floor and redefining background as the complement of candidates at each cutoff. Separately, we retained primary labels but restricted both groups to sizes $\geq10$ and $\geq20$. These are comparisons using existing outputs, not recalibrations under a different testing floor. These sensitivity comparisons did not involve fitting additional regression models.

\subsection{Bayesian hierarchical models}\label{sec:ch6_bayesian_characterisation}

We fitted Bayesian hierarchical models for CSC status and the two screening scores. The results tables use \emph{domain} to label what is analysed: mixing (within-cluster diversity) or composition (sequence-level attributes). Mixing models use continuous entropy values. The tables use \emph{family} to label the model type: logistic for CSC status or linear for mean screening scores. We fitted both model types at both analytical levels. Table~\ref{tab:ch6_model_terminology} defines these labels and the scales used to interpret estimates.

\begin{table}[htbp]
    \centering
    \caption[Terminology for Bayesian characterisation models]{Terminology used in the Bayesian results tables. Domain, model family, predictor scale, and model specification describe separate features of each fitted model.}\label{tab:ch6_model_terminology}
    \begin{thesistablebody}[\thesistablesetup\renewcommand{\arraystretch}{1.5}]{@{}P{0.17\linewidth}P{0.24\linewidth}P{0.53\linewidth}@{}}
    \toprule
    \textbf{Feature} & \textbf{Table label} & \textbf{Meaning and interpretation} \\
    \midrule
    Domain 
        & Mixing & Cluster-level analysis: one row per cluster, with predictors describing diversity of the five attributes. \\
        & Composition & Sequence-level analysis: one row per sequence record, with the five categorical attributes as predictors. \\
    \midrule
    Model family 
        & Logistic & Estimates the probability that the corresponding cluster has CSC status. Coefficients describe differences in log odds; exponentiated coefficients are odds ratios. \\
        & Linear & Estimates mean local burst or onward burden score. Coefficients describe differences in mean score. \\
    \midrule
    Predictor scale 
        & Observed entropy & Each coefficient describes the outcome contrast for a 0.1 increase in normalised entropy, holding other terms constant. \\
        & Null-standardised entropy & Each coefficient describes the outcome contrast for a one-unit increase in $Z^{(H)}$, holding other terms constant. \\
        & Categorical composition & Composition estimates compare each category with its reference group (see Section~\ref{sec:ch6_attributes} above). \\
    \midrule
    Model specification 
        & Primary & Includes the five entropy or categorical predictors jointly, with policy-period and clade varying intercepts. \\
        & Expanded & Adds the three surveillance measures; each surveillance coefficient describes the outcome contrast for a one-SD increase in that measure. \\
    \bottomrule
    \end{thesistablebody}
\end{table}

All models used the same additive structure, without interactions. They estimate an overall association for each attribute, rather than allowing that association to differ with cluster size or between population groups. Cluster size was not included as a separate regression predictor because it is used in the screening process (Chapter~\ref{ch:cluster_transition_screening}). The descriptive analyses across size bands therefore examine a question that these models do not directly answer. The shared model form is

\[
    g(\mu_i)=\alpha+\mathbf{x}_i^{\mathsf T}\boldsymbol{\beta}
    +u_{\mathrm{period}[i]}+v_{\mathrm{clade}[i]},
\]

where $Y_i$ is the outcome and $\mu_i=E(Y_i\mid\mathbf{x}_i,u_{\mathrm{period}[i]},v_{\mathrm{clade}[i]})$ is its conditional expectation. The index $i$ denotes a cluster or sequence record, depending on the analytical level. Here, $\alpha$ is the overall intercept, $\mathbf{x}_i$ contains the attribute predictors and, in expanded models, the surveillance measures included for adjustment, and $\boldsymbol{\beta}$ contains their regression coefficients. The term $\mathbf{x}_i^{\mathsf T}\boldsymbol{\beta}$ sums the predictor values multiplied by their corresponding coefficients. The terms $u$ and $v$ are the policy-period and clade varying intercepts. The link $g$ is logit for Bernoulli models of CSC status and identity for Gaussian models of the screening scores.

The 0.1 increment for observed entropy makes an estimate describe, for example, a change from 0.4 to 0.5 rather than the entire possible range from 0 to 1. The original models implement this by multiplying observed entropy by 10; it is a choice of predictor units, not an additional diversity normalisation. For surveillance measures, one SD refers to the spread of that measure across the relevant cluster or sequence dataset.

All models allow baseline outcomes to vary by policy period and clade through partial pooling: estimates share information across groups, with less precisely estimated baselines drawn towards the overall baseline. These terms describe remaining differences between periods and clades, not causal effects. We fitted 12 mixing models (three outcomes, two entropy scales, and two specifications) and six composition models (three outcomes and two specifications). CSC status and local burst models use size-eligible clusters or their member records; onward burden models use only burden-eligible clusters or their member records. Each model uses all records with complete data for its outcome, predictors, and grouping variables (complete cases), without subsampling or imputation (Appendix~Table~\ref{tab:bayesian_model_specifications}).

In mixing models, attribute coefficients describe associations between within-cluster diversity and the outcome, using the entropy increments in Table~\ref{tab:ch6_model_terminology}; in composition models, they compare each attribute category with its reference group (see Section~\ref{sec:ch6_attributes} above). These associations account for the other predictors and the policy-period and clade intercepts. For CSC status, coefficients describe differences in log odds; exponentiating them gives odds ratios (ORs)\nomenclature{OR}{odds ratio}. Odds are the probability of CSC status divided by the probability of background status, so an odds ratio compares these odds between predictor values. Ratios above or below one indicate higher or lower CSC odds, respectively. For the screening scores, coefficients describe differences in the expected score.

Composition models give larger clusters more weight because they contribute more records. Records within a cluster share the same outcome, and sequences may recur across windows. The models do not include cluster or sequence random intercepts to account for this dependence, so posterior intervals may be too narrow. These models describe the sampled records represented in clusters, rather than individual transmission or infection risks.

Models were fitted using the Bayesian Model Building Interface (Bambi)\nomenclature{Bambi}{BAyesian Model Building Interface}, which uses PyMC \citep{Capretto2022,AbrilPla2023}. We report posterior means, 95\% highest-density intervals (HDIs)\nomenclature{HDI}{highest-density interval}, and the posterior probability of the direction indicated by each estimate. A 95\% HDI is the shortest interval containing 95\% of the posterior probability. Appendix~\ref{sec:app_ch6_bayesian_details} provides further model details, priors, fitting settings, and diagnostic definitions. Sampler diagnostics assess sampling reliability.

\section{Results}\label{sec:ch6_results}

We first examine cluster sizes, then compare population composition and diversity across size bands, screening routes, and descriptive sensitivity comparisons. We subsequently report the Bayesian estimates and their sampling diagnostics.

\subsection{Cluster sizes and the contribution of smaller clusters}\label{sec:ch6_size_results}

The 631 CSC clusters contributed 22,156 sequence records, compared with 130,021 records in 8,331 background clusters (Table~\ref{tab:ch6_size_overview}). CSC clusters had a median size of 22 sequences (IQR 12--37), compared with 9 (7--14) in the background. The distributions nevertheless overlapped, and the background included some very large clusters (Figure~\ref{fig:ch6_size_cumulative}). Raw size alone therefore does not reproduce CSC status, which also depends on context, membership novelty, and onward burden.

\begingroup
    \let\thesistablesetup\thesisdensetablesetup
    \input{assets/scotland/tab_ch6_size_overview}
\endgroup

\begin{figure}[htbp]
    \centering
    \includegraphics[width=\textwidth]{fig_ch6_size_cumulative}
    \caption[Cumulative cluster and record counts by cluster size]{\textbf{How cluster size contributes to the CSC and background comparisons.} For each size on the logarithmic horizontal axis, panels show cumulative clusters (A), their percentage within each group (B), cumulative sequence records in those clusters (C), and their percentage within each group (D). Blue denotes CSC and orange background. All clusters have at least six sequences. Record counts sum cluster memberships across retained windows; a sequence can contribute in more than one window. The contrast between B and D shows why numerous small clusters need not dominate pooled composition.}\label{fig:ch6_size_cumulative}
\end{figure}

Clusters below size 20 accounted for 263 of 631 CSC clusters (41.7\%), but only 2,955 of 22,156 CSC records (13.3\%). In the background, the corresponding proportions were 86.2\% of clusters and 51.9\% of records. Thus, small clusters formed a substantial fraction of CSC clusters but did not dominate the record-weighted CSC composition. The marked difference between cluster-weighted and record-weighted contributions motivates reporting both analytical levels. These curves describe size overlap; they do not measure the screen's ability to distinguish confirmed superspreading from other transmission.

\subsection{Overall population composition and across cluster sizes}\label{sec:ch6_candidate_description}

Overall population composition was similar in the two groups. Males accounted for 48.1\% of CSC records and 47.2\% of background records, a difference of 0.9 percentage points. People aged 15--24 accounted for 19.3\% and 17.6\%, respectively (+1.6 points), and residents of large urban areas for 38.7\% and 38.0\% (+0.7 points). All category differences in the pooled comparison lay between -1.0 and +1.6 percentage points (Figure~\ref{fig:sse_composition_descriptive}). These small absolute differences should be distinguished from relative odds estimated by the regression models. SIMD differences varied by category, with the largest positive difference in quintile 4 (+0.9 points); the categories did not follow a simple ordered pattern.

\begin{figure}[htbp]
    \centering
    \includegraphics[width=\textwidth]{fig_sse_composition_descriptive}
    \caption[Population composition of CSC and background clusters]{\textbf{Overall population composition differs little between CSC and background clusters.} Bars show CSC-minus-background percentage-point differences among 22,156 CSC and 130,021 background sequence records from clusters of size at least six. Each record receives equal weight. Positive (blue) bars indicate higher representation in CSC clusters; negative (red) bars indicate higher representation in background clusters. Panels show sex (A), SIMD quintile (B), age group (C), urban/rural class (D), and health board (E). This is the pooled binary comparison; Figure~\ref{fig:ch6_composition_size} separates cluster sizes.}\label{fig:sse_composition_descriptive}
\end{figure}

The pooled contrasts were not uniform across size bands (Figure~\ref{fig:ch6_composition_size}; Appendix~Table~\ref{tab:app_ch6_size_composition}). Among CSC records, the proportion aged 15--24 was 16.8\%, 16.4\%, 18.6\%, 19.1\%, and 21.6\% across increasing size bands. However, the corresponding CSC-minus-background differences were +2.2, -0.5, -0.9, -0.2, and +1.8 percentage points. The greater representation of young adults in larger CSC clusters therefore did not translate into a steadily increasing contrast with background clusters of similar size. Contrasts for sex and urban/rural class also changed direction: the large-urban contrast ranged from -4.9 points at sizes 50--99 to +3.3 points at sizes $\geq100$. The largest band contained 28 CSC clusters and 98 background clusters, so its record-weighted composition could be influenced by relatively few large clusters.

\begin{figure}[htbp]
    \centering
    \includegraphics[width=\textwidth]{fig_ch6_composition_by_size}
    \caption[Population composition contrasts across cluster sizes]{\textbf{Composition contrasts vary across cluster sizes.} Cells show CSC-minus-background differences in category percentages among pooled records within each size band. Blue indicates higher and red lower representation in CSC clusters; numbers are percentage points. The CSC cluster counts across bands are 105, 158, 264, 76, and 28; background counts are 4,519, 2,664, 884, 166, and 98. Record denominators appear in Table~\ref{tab:ch6_size_overview}; group percentages appear in Appendix~Table~\ref{tab:app_ch6_size_composition}. All displayed attributes are complete. These unadjusted descriptions allow non-monotonic patterns and do not establish size interactions in the regression models.}\label{fig:ch6_composition_size}
\end{figure}

Appendix~Figure~\ref{fig:app_ch6_size_cluster_weighted} repeats these contrasts with equal weight per cluster. Appendix~Figure~\ref{fig:app_composition_profiles} shows the distributions of within-cluster attribute category proportions and their mean contrasts across all eligible clusters.

\subsection{Within-cluster diversity across cluster sizes}\label{sec:ch6_entropy_size_results}

Observed diversity generally increased with cluster size, but a higher observed value did not imply greater diversity relative to the window-specific reference (Figure~\ref{fig:ch6_entropy_size}). For CSC clusters, median observed health-board entropy increased from 0.28 at sizes 6--9 to 0.65 at sizes $\geq100$. Over the same bands, median null-standardised health-board entropy decreased from -2.78 to -7.97. Larger clusters could therefore contain a broader range of residential regions while remaining more concentrated than expected from random samples of that size. Part of the change on the standardised scale may also reflect changes in the reference standard deviation.

\begin{figure}[htbp]
    \centering
    \includegraphics[width=\textwidth,height=0.76\textheight,keepaspectratio]{fig_ch6_entropy_by_size}
    \caption[Observed and null-standardised entropy across cluster sizes]{\textbf{Diversity depends on cluster size on both entropy scales.} Points show medians and bars the interquartile ranges among clusters, separately for CSC (blue) and background (orange). Left: observed entropy on the 0--1 scale. Right: null-standardised entropy, with a dotted line at zero; vertical scales differ by attribute. The IQR describes between-cluster variability, not confidence in a median. All six entropy attributes are complete: CSC counts across bands are 105, 158, 264, 76, and 28, and background counts are 4,519, 2,664, 884, 166, and 98.}\label{fig:ch6_entropy_size}
\end{figure}

The comparison with background depended on size. Overall median null-standardised health-board entropy was -5.00 for CSC and -4.09 for background. Within the smallest band, however, CSC clusters had a less negative median than background (-2.78 versus -3.75), whereas in the largest band the CSC median was more negative (-7.97 versus -5.27). Median null-standardised age entropy in CSC clusters also decreased across the endpoint bands, from -0.17 to -1.45. These patterns show why a size-specific null reference does not remove the need to inspect residual associations with size. They also caution against interpreting a pooled association as applying uniformly to small and large clusters.

\subsection{Screening routes and descriptive sensitivity comparisons}\label{sec:ch6_descriptive_sensitivity_results}

Screening routes selected different size profiles. Burst-only candidates had a median size of 28 sequences, burden-only candidates 9, and both-axis candidates 35.5 (Appendix~Table~\ref{tab:ch6_route_sizes}). The 428 burst-only candidates contributed 18,567 records, or 83.8\% of all CSC records. Thus the pooled CSC composition mainly described the burst-only route. For young adults, the burst-only contrast with all background was +1.7 percentage points, whereas the burden-only contrast with burden-eligible background was +0.2 points (Appendix~Figure~\ref{fig:app_ch6_route_composition}). Median null-standardised health-board entropy was -5.95 for burst-only candidates, compared with -4.09 for all background, but -3.12 for burden-only candidates, compared with -4.32 for burden-eligible background (Appendix~Figure~\ref{fig:app_ch6_route_entropy}). These descriptive contrasts are consistent with different population profiles across screening routes, while leaving size, timing, and other contextual differences unresolved.

Changing the cutoff also changed composition. At $p\leq0.01$, 129 clusters were selected and the young-adult contrast was +0.1 percentage points, compared with +1.6 points under the primary cutoff and +2.2 points at $p\leq0.10$ (1,230 clusters). The large-urban contrast changed from +0.7 points in the primary comparison to -2.8 points at $p\leq0.01$. A stricter cutoff therefore did not simply amplify the primary demographic profile (Appendix~Table~\ref{tab:ch6_sensitivity_sizes} and Figure~\ref{fig:app_ch6_sensitivity_composition}).

Keeping primary labels but restricting both groups to sizes $\geq20$ retained 368 CSC and 1,148 background clusters. Their young-adult proportions were 19.7\% and 19.6\%, a difference of only +0.1 percentage points. This attenuation suggests that different size distributions contribute to the pooled contrast; it does not establish that small CSC clusters were misclassified. Entropy summaries under all cutoff and size restrictions are shown in Appendix~Figure~\ref{fig:app_ch6_sensitivity_entropy}. None of these descriptive comparisons validates the screen against confirmed events or replaces adjustment for differences in epidemic context.

\subsection{Bayesian characterisation results}\label{sec:ch6_bayesian_results}

Unless stated otherwise, estimates below are from the expanded models that adjusted for the three surveillance measures (Section~\ref{sec:ch6_attributes}), and are reported as posterior means with 95\% HDIs. Estimates for the five demographic and geographic attributes were generally similar between primary and expanded models, particularly for CSC status. Larger differences occurred for some estimates from the composition models for the screening scores. Appendix~Table~\ref{tab:bayesian_fixed_effects_main} summarises the focal attribute estimates. Complete coefficient estimates for all specifications are provided as electronic supplement (Appendix~\ref{sec:app_ch6_bayesian_results}).

\subsubsection*{Model diagnostics}\label{sec:ch6_model_diagnostics}

All 18 models had zero divergent transitions, and 15 met the specified sampling criteria (Appendix~Table~\ref{tab:bayesian_model_diagnostics}). Warnings concerned $\widehat R$ and/or bulk effective sample size (ESS)\nomenclature{ESS (sampling)}{effective sample size; Bayesian sampling diagnostic} in three sequence-level composition models: both local burst models and the expanded onward burden model. All cluster-level mixing models and both composition CSC models met the criteria. Diagnostic status is retained alongside the estimates in the appendix tables; exact fitted sample sizes are reported in Appendix~Table~\ref{tab:bayesian_model_specifications}.

\subsubsection*{Associations with within-cluster diversity}

Greater age diversity relative to the simulated reference (higher null-standardised entropy) was associated with lower CSC odds (OR 0.85, 95\% HDI 0.79--0.91 per one-unit increase), whereas greater diversity in urban/rural class on the same scale was associated with higher CSC odds (OR 1.07, 1.01--1.13; Figure~\ref{fig:fixed_effects_mixing}A). Sex and health-board entropy also showed negative associations, with a smaller negative estimate for SIMD. Local burst associations followed the same directions (Figure~\ref{fig:fixed_effects_mixing}B); the largest negative coefficient was for health-board entropy ($\beta$ -0.030, 95\% HDI -0.032 to -0.028). Onward burden associations differed (Figure~\ref{fig:fixed_effects_mixing}C): health-board and SIMD entropy were positively associated with score, entropy for urban/rural class was negatively associated, and age and sex had unclear associations (HDIs crossed zero).

\begin{figure}[htbp]
    \centering
    \includegraphics[width=\textwidth]{fig_fixed_effects_mixing}
    \caption[Bayesian fixed-effect estimates for mixing models]{
        \textbf{Bayesian fixed-effect estimates for cluster-level mixing models.} Points show posterior means and horizontal bars show 95\% HDIs for primary and expanded models on both entropy scales, distinguished by colour and symbol. Panel A shows CSC odds ratios; Panels B and C show differences in mean local burst and onward burden scores. Estimates correspond to a 0.1 increase in observed entropy or a one-unit increase in null-standardised entropy. Dashed lines mark no association (one for odds ratios; zero for score differences). Both specifications include the five entropy predictors jointly and policy-period and clade varying intercepts; expanded models additionally include the three standardised surveillance measures. All models shown met the sampling criteria.}\label{fig:fixed_effects_mixing}
\end{figure}

On the observed scale, entropy for age, SIMD, health board, and urban/rural class was positively associated with CSC status (ORs 1.08--1.29 per 0.1 increase in observed entropy); the HDI for sex spanned one. All five observed-entropy coefficients were positive for local burst score. For onward burden score, observed entropy for sex, age, SIMD, and urban/rural class was negatively associated, while health-board entropy was positively associated.

\subsubsection*{Associations with sequence-level composition}

Figure~\ref{fig:fixed_effects_composition} reports the adjusted composition contrasts. Records from females had lower odds of belonging to a CSC cluster and lower mean scores than records from males. Records from people aged 15--24 had higher CSC odds than those from people aged 25--64 (OR 1.11, 95\% HDI 1.06--1.15) and higher mean scores; children and older adults generally showed the opposite pattern. Remote towns and remote rural areas had lower CSC odds and lower mean scores than large urban areas; the CSC odds ratio for remote towns was 0.80 (0.71--0.88). SIMD contrasts varied by category without a simple ordered pattern; the categorical specification did not impose a monotonic trend.

\begin{figure}[htbp]
    \centering
    \includegraphics[width=\textwidth]{fig_fixed_effects_composition}
    \caption[Bayesian fixed-effect estimates for composition models]{
        \textbf{Bayesian fixed-effect estimates for sequence-level composition models.} Points show posterior means and horizontal bars show 95\% HDIs for expanded (blue squares) and primary (orange circles) models. Panels A and B show CSC odds ratios; C and D show mean local burst score differences; E and F show mean onward burden score differences. Contrasts are relative to male sex, age 25--64, SIMD quintile 1, large urban areas, and Greater Glasgow and Clyde. Dashed lines mark one for odds ratios and zero for score differences. Both specifications jointly include the five attributes and policy-period and clade varying intercepts; expanded models add the three standardised surveillance measures. Orkney and Western Isles are omitted from B and F for readability; their estimates appear in Appendix~Table~\ref{tab:bayesian_fixed_effects_main}. Both local burst models and the expanded onward burden model had sampling warnings (Appendix~Table~\ref{tab:bayesian_model_diagnostics}).}\label{fig:fixed_effects_composition}
\end{figure}

Differences between health boards varied more in size and direction. Relative to Greater Glasgow and Clyde, Orkney had the largest positive CSC estimate (OR 3.82, 95\% HDI 3.00--4.95); Forth Valley was also higher, while several mainland boards had lower estimates. Geographic differences varied across outcomes: Orkney had a higher local burst score but a lower onward burden score. Both expanded composition models for screening scores had the sampling warnings described above.

\subsubsection*{Surveillance adjustment and model comparison}\label{sec:ch6_primary_expanded_adjustment}

Adding the three surveillance measures generally changed the estimated associations of sex, age, SIMD, urban/rural class, and health board only slightly (Figures~\ref{fig:fixed_effects_mixing} and~\ref{fig:fixed_effects_composition}). In mixing models, all five CSC odds ratios based on null-standardised entropy were unchanged to two decimal places. The reported CSC odds ratios based on observed entropy changed by at most 0.02, while coefficients for both screening scores changed by no more than 0.002 on either entropy scale.

Composition estimates for CSC status were also generally similar, although some health-board associations became less certain. For Highland relative to Greater Glasgow and Clyde, the OR decreased from 1.12 (95\% HDI 1.03--1.22) in the primary model to 1.07 (0.98--1.16) in the expanded model. Larger changes occurred in composition estimates for the screening scores, particularly for geographic contrasts: the onward burden coefficient for remote rural areas relative to large urban areas changed from -0.012 (95\% HDI -0.024 to approximately zero) to -0.027 (-0.039 to -0.015). 

Surveillance associations appear in Appendix~Figure~\ref{fig:app_random_effects_mixing} for mixing models and Figure~\ref{fig:random_effects_composition} for composition models. The Chapter~6 electronic supplement provides the corresponding numerical estimates for primary and expanded models (Appendix~\ref{sec:app_ch6_bayesian_results}).

\begin{figure}[htbp]
    \centering
    \includegraphics[width=\textwidth]{fig_random_effects_composition}
    \caption[Bayesian surveillance and varying-intercept estimates for composition models]{\textbf{Bayesian surveillance and varying-intercept estimates for sequence-level composition models.} Points show posterior means and horizontal bars show 95\% HDIs for expanded (blue squares) and primary (orange circles) models. The first three rows show surveillance coefficients, included only in expanded models, per one-SD increase. Remaining rows show differences in the model baseline between policy periods and between clades. Panel A uses a logarithmic axis: surveillance estimates are odds ratios, while period and clade estimates are odds multipliers relative to the model baseline, with other model terms held constant. Panels B and C show differences in mean local burst and onward burden scores. Dashed lines mark one in A and zero in B and C. Policy-period codes are defined in Chapter~\ref{ch:genomic_networks}. Both local burst models and the expanded onward burden model had sampling warnings (Appendix~Table~\ref{tab:bayesian_model_diagnostics}).}\label{fig:random_effects_composition}
\end{figure}

\subsubsection*{Policy-period and clade variation}\label{sec:ch6_period_clade_results}

For CSC status, period and clade estimates are reported as odds multipliers relative to a zero value of the corresponding varying-intercept term, holding all other model terms constant. A multiplier of two therefore means twice the CSC odds. In the expanded composition CSC model, the highest period odds multipliers occurred during route map phase~3 (P3), pre-tier tightening (T1), and the five-tier framework (F5) (Figure~\ref{fig:random_effects_composition}). T1 had the largest posterior mean odds multiplier (9.49, 95\% HDI 4.48--20.09); final easing (FE) and the post-restriction period (PR) had the lowest estimates. Period odds multipliers in the mixing CSC models were closer to one, with intervals generally spanning one (Appendix~Figure~\ref{fig:app_random_effects_mixing}).

In the expanded composition CSC model, clades 20B and 20E had odds multipliers below one, while 22D had the largest posterior mean multiplier, although with substantial uncertainty (42.52, 95\% HDI 13.46--148.41). Several other Omicron-era clades also had multipliers above one, while others had multipliers below one or intervals spanning one. Across all three expanded composition models, the posterior mean group-level standard deviations indicated greater variation between clades than between policy periods (Appendix~Table~\ref{tab:app_bayesian_random_effect_sds}).

\section{Discussion}\label{sec:ch6_discussion}

This chapter asked how population composition and within-cluster diversity were associated with CSC status and the two screening scores, and whether local bursts and onward burden showed different population profiles.

The main contribution is a description of how the populations represented by the genomic screen vary with cluster size and screening route. Pooled CSC and background composition differed by at most 1.6 percentage points, despite different size distributions. The contrasts for age and urban/rural class changed across size bands and screening cutoffs, while burst-only candidates contributed most CSC records. Consequently, the pooled CSC profile is not a single population signature shared equally by every type or size of candidate.

Entropy adds a complementary description: clusters may span several demographic groups or regions while remaining more concentrated than random samples of comparable size from the same window. The opposite health-board associations for local burst and onward burden scores suggest hypotheses about concentrated exposure and wider connections, consistent with the differing consequences of documented outbreaks \citep{Lemieux2021}. These interpretations concern population structure within the sampled record. They cannot identify a setting, a transmission route, or a causal risk factor without additional epidemiological evidence.

\subsection{Demographic and geographic profiles}\label{sec:ch6_discussion_composition}

Young adults aged 15--24 were more represented in CSC clusters than adults aged 25--64, even after adjustment for the other measured attributes, policy period, viral clade, and surveillance measures. The adjusted comparison concerned relative odds: sampled records from this age group had 11\% higher odds of belonging to a CSC cluster (OR 1.11, 95\% HDI 1.06--1.15). The pooled absolute young-adult difference was 1.6 percentage points and fell to 0.1 points when both groups were restricted to clusters of size at least 20. These summaries address different comparisons: the regression adjusted for attributes, policy period, clade, and surveillance measures, but did not include cluster size or its interactions. Its age coefficient therefore cannot establish that young-adult representation differs consistently between CSC and background at every cluster size. These comparisons describe membership of screened clusters among sequenced infections; they do not estimate an individual's risk of becoming infected or transmitting the virus.

Differences in contact patterns are one possible explanation for population concentration. Age-assortative contacts documented in European contact surveys provide a rationale for examining age diversity as well as age composition \citep{Mossong2008}. Contact tracing research found that transmission risk increased with longer exposure and closer social contact \citep{Sun2021}. In Scotland, the EURO 2020 investigation also found that cases reporting attendance at football-related events were predominantly young adult men \citep{Marsh2021}. These studies show how particular social activities can produce a distinct demographic profile. The present data do not contain the exposure histories needed to identify which activities, if any, explain the CSC profiles.

Deprivation provided less distinction between CSC and background clusters. The categorical SIMD contrasts did not follow a simple ordered pattern, and some varied with cluster size. This finding addresses a different question from studies reporting higher infection risk in more deprived populations \citep{Niedzwiedz2020}: here, the comparison was between sequenced infections in candidate and background clusters. The absence of a clear gradient therefore does not show that deprivation was unimportant to the wider epidemic.

Health-board differences were more variable than the demographic contrasts. Orkney, for example, had higher CSC odds than Greater Glasgow and Clyde, but its estimated local burst score was higher and its onward burden score lower. This suggests a profile worth examining in local outbreak records, although the score models retained sampling warnings. Local introductions, outbreak investigations, and sequencing practices could all contribute to such geographic differences. The estimates cannot establish that a place has an inherently greater risk of superspreading.

\subsection{Population composition and within-cluster diversity}\label{sec:ch6_discussion_mixing}

A cluster can contain people from several population groups and still be concentrated in one of them. For example, a large cluster may span several age groups while most of its members belong to a single age group. This helps explain why the two diversity measures gave different results. Observed entropy describes how mixed a cluster is in absolute terms. Null-standardised entropy compares that diversity with random samples of the same size from the same surveillance window. The latter supplies an explicit reference for the greater opportunity for large clusters to include several groups and for changes in the sampled population over time. Nevertheless, the size-band results show that standardised values retain associations with size. The changing reference standard deviation and the underlying composition of larger clusters can both contribute.

On the null-standardised scale, CSC status was associated with lower diversity in age, sex, and health board, with a weaker association for deprivation. This suggests that demographic and regional concentration is relevant to the patterns selected by the screen. The geographic result complements Chapter~\ref{ch:genomic_networks}, where higher pairwise compatibility scores tended to occur among infections from the same residential area. These are overall associations across a continuous diversity scale; they do not establish that every CSC cluster was less diverse than its simulated reference. CSC clusters had varied diversity profiles, and the direction of the health-board contrast with background was not the same in every size band. A cluster with relatively concentrated attributes could arise from shared exposure or repeated transmission among similar people. A broader distribution could arise from an exposure involving several groups or from subsequent dissemination. Entropy alone cannot distinguish these processes.

Diversity in urban/rural class showed the opposite adjusted association: clusters containing a greater mix of urban/rural classes had higher CSC odds. Regional concentration and mixing across urban/rural classes can coexist, because one health board can include urban and rural communities. Commuting or shared workplaces and institutions could produce such a pattern, but these are explanations to investigate. Early Scottish genomic research documented community, healthcare, and conference outbreaks following viral introductions \citep{daSilvaFilipe2020}; residential categories alone cannot distinguish those exposure settings.

\subsection{Local bursts and onward burden}\label{sec:ch6_discussion_burst_burden_context}

Geography showed the clearest contrast between the two screening scores. Lower health-board diversity on the null-standardised scale was associated with higher local burst scores, whereas greater health-board diversity was associated with higher onward burden scores. Greater deprivation diversity also accompanied higher onward burden, while age and sex diversity had no clear associations with this score. Thus, the population profile associated with an unusually large or novel cluster differed from the profile associated with more sampled infections in later linked clusters relative to the starting cluster's size.

One hypothesis is that local bursts reflect concentrated transmission within a population, while connections across regions help sustain wider spread. Documented outbreaks show why this distinction matters. In Boston, a nursing-facility outbreak involved intense local transmission with little wider spread, whereas a conference outbreak was followed by widespread transmission \citep{Lemieux2021}. This example suggests testing whether different contact patterns explain the two Scottish screening signals. The transition graph itself links clusters that share sampled sequences and cannot establish transmission from one cluster to another.

The screen's design could also contribute to these differences. The local burst score includes cluster size. Onward burden is measured relative to the starting cluster's size and assessed only for clusters with positive observed downstream burden. Moreover, 428 candidates were burst-only, 191 were burden-only, and just 12 met both criteria (Chapter~\ref{ch:cluster_transition_screening}). The descriptive route comparison reinforces this distinction: burst-only candidates contributed 83.8\% of CSC records and had lower median null-standardised health-board diversity than all background, whereas burden-only candidates had higher median diversity than burden-eligible background. This helps explain why the combined profile resembles the local burst profile. However, burden-only candidates were much smaller than burst-only candidates. Comparing these routes within a model that represents size and context would be a further sensitivity analysis, rather than a conclusion already established by the current models.

\subsection{Epidemic timing and surveillance}\label{sec:ch6_discussion_context}

CSC profiles also varied with epidemic period and viral clade. In the composition model, higher period estimates occurred during the 2020 reopening phase (P3), subsequent tightening of restrictions (T1), and introduction of the five-tier framework (F5). Clade estimates also varied widely. Such differences may reflect when and where viruses were introduced, opportunities for contact, population immunity, and which infections were detected. UK genomic research similarly found that earlier transmission lineages were larger and more geographically dispersed \citep{duPlessis2021}. Because policy periods and clades overlap in time, the present models cannot separate policy effects from viral characteristics and other changes in the epidemic.

Adding the three incidence and sequencing measures changed most CSC associations little, giving some reassurance that these measured differences did not account for the main profiles. Some score associations changed more. Higher window sequencing proportions remained associated with higher CSC odds in the composition model. These adjustments cannot recover untested infections or fully account for targeted outbreak sequencing. Simulations by \citet{Han2023} likewise showed that low testing rates and uneven sampling could delay variant detection even when a larger proportion of available samples was sequenced.

\subsection{Interpretation and next steps}\label{sec:ch6_discussion_implications}

A strength of this analysis is that demographic and geographic attributes were examined after screening rather than used to assign CSC status. The additional size analyses address a specific concern about the binary comparison: the CSC record pool was not dominated by clusters below size 20, but differences in size distributions contributed to the pooled demographic contrasts. Small absolute differences may also reflect genuinely similar sampled populations, heterogeneity between screening routes, or limited discrimination by the screen. Without independently labelled outbreaks, the present data cannot separate these explanations or estimate screening accuracy. A stricter cutoff did not consistently strengthen the composition contrasts, so greater statistical extremeness of a score should not be equated with a stronger demographic signature.

The models estimate additive associations across their fitted datasets. They do not estimate interactions between size and attributes, and their coefficients should not be described as associations independent of cluster size. The diversity reference conditions on size and window composition but does not preserve lineage structure or relationships between attributes. Further cluster-level analyses could examine size flexibly and allow selected interactions, particularly for age and geographic diversity. Since cluster size is itself part of the local burst screen, such adjustment would estimate a conditional association with the screen, not a causal effect of population characteristics on transmission.

The composition estimates also require care. Larger clusters contribute more records, members of the same cluster share an outcome, and a sequence can appear in more than one window. The models do not fully account for this dependence, so their uncertainty intervals may be too narrow. Three composition models for screening scores retained sampling warnings, and posterior predictive checks were not reported. Future model development should assess these issues together with size dependence. The descriptive size analyses help assess heterogeneity in the sampled populations, but do not resolve these modelling limitations.

Two hypotheses can guide the next stage. First, candidates with unusually low age diversity may more often correspond to documented shared activities or institutions than clusters of similar size from the same epidemic context. Linking to exposure and contact tracing records would allow comparison of the frequency of shared settings; similar frequencies across these groups would weaken this explanation.

Second, greater health-board diversity among clusters with onward burden may correspond to epidemiologically documented links across communities. This could be tested using locations and timing of linked cases, comparing clusters with similar size, surveillance coverage, and opportunity for follow-up. An association that disappears after accounting for introductions or targeted sampling would favour an observation-process explanation.

These are hypotheses for external investigation, not evidence that residential diversity measures contact between communities. Their value is to identify what additional information is needed to explain the contrasting screening profiles. The small and sometimes unstable composition contrasts do not, by themselves, justify targeting particular demographic groups or places as sources of superspreading.

\section{Conclusion}\label{sec:ch6_discussion_conclusions}

CSC clusters were larger than background clusters, but their pooled demographic and geographic composition was similar, with category differences of at most 1.6 percentage points. Smaller clusters did not dominate CSC sequence records, yet composition contrasts varied across size bands and screening cutoffs. Burst-only candidates contributed most CSC records, while burden-only candidates had a different size and geographic-diversity profile.

Entropy described population concentration and breadth beyond the identity of the represented categories. Both observed and null-standardised entropy varied with cluster size, so the regression associations should not be interpreted as independent of size. The contrast between regional concentration in local bursts and broader geographic representation with onward burden motivates specific tests using outbreak and contact tracing records. Those links are needed to determine whether the profiles reflect shared exposures, wider dissemination, or surveillance processes.
